\section{Inference and Evaluation Methodology}

Anomaly detection models on time series data face a fundamental discrepancy between how they are \textit{trained} and how their real-world performance is \textit{evaluated}. This project employs two distinct evaluation paradigms: Window-Level Evaluation (used strictly during training) and Event-Level Evaluation (used during inference and robustness testing).

\subsection{Window-Level Evaluation (Training Script)}
During training (e.g., in \texttt{train.py}), the continuous time series is sliced into thousands of overlapping windows of a fixed size (e.g., 512 points). A window is labeled as an anomaly if anomalous points occupy more than a strict threshold (e.g., $10\%$) of the window's span.

\textbf{F1 Score Penalty:} 
Because the dataset relies on sliding windows, many generated windows capture only the extreme edges of an anomaly (e.g., $85\%$ background noise and $15\%$ anomaly curve). These "boundary" windows are highly ambiguous and often misclassified by the network. During training, the Macro F1 score evaluates each window independently. Consequently, the F1 score appears artificially deflated because the network is heavily penalized for hesitating on these blurry edges, especially under severe class imbalance conditions.

\subsection{Event-Level Evaluation (Inference Script)}
To provide a rigorous assessment of practical detection performance, the inference framework evaluates the model dynamically using detection-driven event segmentation, decoupling the evaluation from rigid window boundaries.

\textbf{1. Probability Accumulation Over Overlapping Windows} \\
During inference, the continuous univariate time series is partitioned into overlapping windows of a fixed length (e.g., 512 points) moving forward by a smaller stride (e.g., 10 points). 

For each window, the model outputs a probability distribution over the available classes via a softmax layer. Because the stride is much smaller than the window length, any arbitrary time step is covered by multiple overlapping windows. The final point-wise probability for each time step is calculated as the arithmetic mean of the probabilities predicted by all windows that overlap that specific point. This aggregation acts as a temporal smoothing filter, mitigating classification instability at anomaly boundaries.

\textbf{2. Detection-Driven Event Segmentation} \\
In a streaming context, the boundaries of an anomalous event are unknown \textit{a priori}. Therefore, the framework defines a predicted event entirely via the model's point-wise confidence for the normal background class. 

The evaluation script first constructs a binary anomaly mask, flagging any point where the probability of the background class drops below a strict threshold (e.g., $30\%$). To prevent signal fragmentation caused by transient stochastic noise or instantaneous probability spikes, a temporal morphological closing operator is applied. Specifically, if two anomalous intervals are separated by a small gap (e.g., under 50 points), the intervening frames are bridged, treating them as a single cohesive event. 

A predicted event is then formalized as a maximal contiguous interval of flagged points.

\textbf{3. Event Classification and Scoring} \\
Once the temporal boundaries of a predicted event are established, the framework must determine the event's overall classification. Rather than evaluating points individually, the aggregated probability for each anomaly class is integrated (summed) over the entire duration of the predicted event. The predicted event is subsequently assigned the classification of the anomaly class with the largest area under the curve.

The final evaluation metric calculates the temporal overlap between these predicted events and the ground-truth anomaly events. A True Positive is recorded if a predicted event intersects with a ground-truth event of the corresponding class. A False Positive is recorded if the model predicts an event in pure background noise, or assigns the wrong classification. A False Negative is recorded if a ground-truth event has no overlapping prediction. This logic fully decouples the evaluation from rigid point-wise constraints, yielding a metric that accurately reflects real-world operational reliability.
